\section{Homología simplicial}

Hasta ahora hemos demostrado que dada una variedad diferenciable de dimensión, podemos, mediante una función de Morse, extraer su estructura de CW-complejo. Es decir, reconstruirla a partir del pegado de celdas cuyas dimensiones vienen determinadas por el índice de cada uno de los puntos críticos de una función de Morse.

De cara al siguiente y último capítulo sobre las \textit{desigualdades de Morse} necesitamos introducir los \textit{grupos de homología} de una variedad y sus correspondientes \textit{números de Betti}. Desarrollaremos estos conceptos en el contexto de los espacios topológicos \textit{triangulados}.

Cuando queremos estudiar espacios topológicos de dimensión baja, el grupo fundamental $\pi_1(X)$ resulta muy útil. Sin embargo, este no nos permite diferenciar entre espacios de dimensiones más altas, por ejemplo el grupo fundamental de $\mathbb{S}^n$ es trivial para $n\ge 2$. Esto puede solucionarse considerando los grupos de homotopía $\pi_n(X)$, el problema es que estos pueden resultar extremadamente complicados de computar. Esta idea queda reforzada cuando uno se da cuenta de que ni siquiera existe una manera robusta de computar los grupos de homotopía de $\mathbb{S}^n$ para cualquier $n>1$. Ponemos como ejemplo algunos de los grupos de homotopía conocidos de $\mathbb{S}^2$ para tener una idea de su complejidad:
\begin{center}
    \begin{tabular}{c|ccccccccccccccc}
    & $\pi_1$ & $\pi_2$ & $\pi_3$ & $\pi_4$ & $\pi_5$ & $\pi_6$ & $\pi_7$ & $\pi_8$ & $\pi_9$ & $\pi_{10}$ & $\pi_{11}$ & $\pi_{12}$ & $\pi_{13}$ & $\pi_{14}$ & $\pi_{15}$ \\
    \hline
    $\mathbb{S}^2$ & $0$ & $\mathbb{Z}$ & $\mathbb{Z}$ & $\mathbb{Z}_2$ & $\mathbb{Z}_2$ & $\mathbb{Z}_{12}$ & $\mathbb{Z}_2$ & $\mathbb{Z}_2$ & $\mathbb{Z}_3$ & $\mathbb{Z}_{15}$ & $\mathbb{Z}_2$ & $\mathbb{Z}_2^2$ & $\mathbb{Z}_{12}\times \mathbb{Z}_2$ & $\mathbb{Z}_{84}\times \mathbb{Z}_2^2$ & $\mathbb{Z}_2^2$
    \end{tabular}
\end{center}
\vspace{1em}

Por suerte, existen los grupos de homología $H_n(X)$, estos son más sencillos de computar y también resuelven limitaciones del grupo fundamental al tratar con espacios de dimensiones altas. El precio a pagar por una mayor computabilidad se traduce en tener que desarrolllar una teoría que no solo es compleja si no que a priori resulta más opaca.

De nuevo, estamos de suerte, porque para la tarea que nos ocupa nos basta con la llamada \textit{homología simplicial}. Esta es una versión restringida de la más complicada y general \textit{homología singular} y, aunque habrá aspectos de esta que debamos introducir, el vocabulario que necesitamos es definitivamente más sencillo. La línea general que seguiremos para el desarrollo de esta teoría es la que se utiliza en \cite{munoz2021}, partiendo de un espacio topológico triangulado.


\bigskip

%%%%%%%%%%%%%%%%%%%%%%%%%%%%%%%%%%%%%%%%%%%%%%%%
%%                                            %%
%%             TRIANGULACIONES                %%
%%                                            %%
%%%%%%%%%%%%%%%%%%%%%%%%%%%%%%%%%%%%%%%%%%%%%%%%
\subsection{Triangulaciones}
\label{subsec:triangulaciones}

Triangular un espacio topológico consiste esencialmente en expresarlo como la unión de \textit{poliedros} de su misma dimensión adecuadamente pegados.

\begin{definition} \hfill

    \begin{enumerate}[label=(\roman*)]
        \item  Un \textit{k-poliedro convexo}, $P^k$, es la envoltura convexa de una cantidad finita de puntos de $\mathbb{R}^n$ con interior no vacío.
        \item Una \textit{l-cara} de $P^k$ es un subconjunto $P^l\subset \partial P^k$ obtenido como la intersección de $P^k$ con un hiperplano que no corta al su interior y que tiene dimensión $l$
    \end{enumerate}
   
\end{definition}

\bigskip

\begin{definition}
    Una \textit{triangulación} de un espacio topológico $X$ se denota 
    \begin{equation*}
        \tau = \{(P^k_{\alpha}, f_{\alpha}): ~0\le k \le n, ~\alpha\in \Lambda_k\}
    \end{equation*}
    y consiste de:
    \begin{enumerate}
        \item Una familia de poliedros $P^k_{\alpha}$ de dimensiones entre $0$ y $n$.
        \item Para cada poliedro, una función continua $f_{\alpha}:P^k_{\alpha}\to X$ con las siguientes propiedades:
        \begin{enumerate}
            \item $f_{\alpha}: \mathring{P}^k_{\alpha}\to f_{\alpha}(\mathring{P}^k_{\alpha})\subset X$ es un homeomorfismo.
            \item Los conjuntos $C^k_{\alpha}=f_{\alpha}(P^k_{\alpha})\subset X$ forman un recubrimiento por cerrados localmente finito de $X$.
            \item Las imágenes $f_{\alpha}(\mathring{P}^k_{\alpha})\subset X$ fromar un recubrimiento disjunto de $X$.
            \item Para cada $l$-cara $P^l\in \partial P^k_{\alpha}$ existe una función en la triangulación $f_{\beta}:P^l_{\beta}\to X$ tal que $f_{\beta} = f_{\alpha}\circ L$, donde $L: P^l_{\beta}\to P^l\subset P^k_{\alpha}$ es un isomorfismo afín seguido de la inclusión.
        \end{enumerate}
    \end{enumerate}
\end{definition}

\bigskip

Cada $C^k_{\alpha}$ es una $k$-celda y $n=dim(X)$ la dimensión máxima de los poliedros de la triangulación.

Recordamos que un recubrimiento $\mathcal{C}=\{C_i\}$ de un espacio topológico $X$ se dice \textit{localmente finito} si todo $x\in X$ tiene un entorno abierto $\mathcal{U}_x$ que interseca solamente a una cantidad finita de elementos de $\mathcal{C}$. 

Diremos que la topología de un espacio $X$ es \textit{coherente} respecto de un recubrimiento $\mathcal{C}$ si la intersección de todo abierto $\mathcal{U}$ de $X$ con cualquier cerrado $C_i$ del recubrimiento es abierta en $C_i$.

Se tiene que, dado un recubrimiento por cerrados localmente finito $\mathcal{C}$ de un espacio $X$, la topología de $X$ es coherente respecto de $\mathcal{C}$. Es decir, la topología de un espacio topológico $X$ es coherente respecto del recubrimiento obtenido al triangular. Esto nos permite reconstruir la topología de $X$ teniendo solo una triangulación suya como la coleción de subconjuntos $\mathcal{U}\subseteq X$ tales que $\mathcal{U}\cap C^k_{\alpha}$ es abierto en $C^k_{\alpha}$ para cualquier $\alpha\in \Lambda_k$.

La propiedad $2.d$ determina que las intersecciones de imágenes de dos poliedros $P^k_{\alpha}$ y $P^t_{\beta}$ de una triangulación están formadas por $l$-celdas, no pueden intersecar sus interiores. Además, a cada función $f_{\alpha}$ le permitimos mandar más de una $l$-cara a una misma $l$-celda.


\bigskip

\begin{definition}
    Con la notación de la definición anterior, una \textit{triangulación regular} es una triangulación tal que
    \begin{enumerate}
        \item $f_{\alpha}: P^k_{\alpha} \to C^k_{\alpha}$ es un homeomorfismo para todo $\alpha\in\Lambda_k$.
        \item Los $P^k_{\alpha}$ son \textit{k-símplices}, es decir, la envoltura convexa de $(k+1)$ puntos afinmente independientes.
        \item Cualesquiera dos $k$-celdas no pueden compartir todos sus vértices.
    \end{enumerate}
\end{definition}

\begin{remark}
    Un $k$-símplice es el objeto geométrico más simple posible de dimensión $k$, el análogo $k$-dimensional del triángulo. Todo $k$-símplice es un $k$-poliedro pero no al revés. De hecho, todo $k$-poliedro puede obtenerse pegando adecuadamente $k$-símplices.
\end{remark}

\bigskip

Una triangulación regular resulta más sencilla de visualizar. En una triangulación regular dos $k$-símplices distintos cualesquiera o no se tocan o comparten exactamente un único $l$-símplice (con $l<k$) de sus fronteras. Representa lo que uno ve al pensar en una malla de triángulos, en el caso de dimensión $2$. 

Por otro lado, al llevar sobre un espacio los poliedros de una triangulación no regular, pueden aparecer identificaciones de las caras de un mismo poliedro. En dimensión $2$, un triángulo puede convertirse por ejemplo en un ``cono'' al quedar dos de sus aristas pegadas mediante la función que lo lleva al espacio que triangula. Resultando en dos posibles inclusiones del vértice que no era común a ambas en dicho espacio.

\bigskip

Finalmente, introducimos la noción de orientación de un poliedro.

\begin{definition}
    Una \textit{orientación} $o$ de un $k$-poliedro es una orientación de su interior $\mathring{P}^k$, es decir, una orientación de $\mathbb{R}^n$.
\end{definition}

Como $\mathring{P}^k$ es un abierto de $\mathbb{R}^k$, su orientación se obtiene fijando una de los dos orientaciones posibles de $\mathbb{R}^k$. Esta orientación $o$ induce una orientación $o|_{P^{n-1}}$ en cada una de las $(n-1)$-caras $P^{n-1}\subset \partial P^k$.

Más concretamente, una base de $\mathbb{R}^k$ pertenece a una de dos clases de equivalencia de bases bajo la relación siguiente: dos bases $B_1$ y $B_2$ están relacionadas si el determinante de la matriz de cambio de base es positivo. Orientar $\mathbb{R}^k$ es sencillamente elegir una de estas dos clases.

Por otro lado, dado un $k$-poliedro orientado $P^k$, la orientación inducida sobre un $k-1$-poliedro $P^{k-1}$ de su frontera es de nuevo una elección de clase de equivalencia de bases de $\mathbb{R}^{k-1}$ pero esta vez forzada por la que la induce. Supongamos que la orientación del $P^k$ es la dada por la clase $[v_1,v_2,\dots, v_k]$, entonces la orientación inducidad sobre $P^{k-1}$ será la dada por la clase $[w_2,\dots,w_k]$ que verifica que
\begin{equation*}
    [\nu, w_2,\dots, w_k] = [v_1,\dots,v_k]
\end{equation*}
donde $\nu$ es un vector tangente a $P^k$, ortogonal a $P^{k-1}$ y apunta hacia el exterior de $P^k$.

En el caso $0$-dimensional, es decir, los vértices, la orientación será la elección del signo y utilizaremos la siguiente convención: dados una arista $a$ y un vértice $v$, $a$ induce la orientación $o(v)=\{+\}$ si termina en $v$ y $o(v)=\{-\}$ si empieza en $v$. Si $a=[v_1,v_2]$ y el orden determina la orientación $a$, entonces $o(v_1)=\{-\}$ y $o(v_2)=\{+\}$.

Es conveniente señalar que, dada una triangulación regular, si dos $k$-poliedros con la misma orientación comparten un $k-1$-poliedro de sus fronteras entonces inducen orientaciones opuestas en el mismo. Basta observar que los vectores $\nu_1$ y $\nu_2$ correspondientes son opuestos.

Existen otras nociones triangulación, como la \textit{pseudotriangulación} y toda una teoría de las llamadas \textit{variedades PL} que para nuestro propósito no son necesarias. Para mayor profundidad en este tema, detalles técnicos y ejemplos sugerimos que el lector se refiera a la fuente original \cite{munoz2021}.

\bigskip

Terminamos la sección de triangulaciones con un resultado de Whitehead recogido en \cite{Whitehead1978} que nos permitirá hablar de la homología de cualquier variedad diferenciable.

\begin{theorem}
    Para toda variedad diferenciable $n$-dimensional $M$ existe una triangulación $\tau$.
\end{theorem}

Este resultado se corresponde con el teorema $7$, situado en la página $819$ de \cite{Whitehead1978}. Primero Whitehead lo formula para variedades cerradas y más tarde, en el apartado $4$ (página $823$) lo generaliza para las variedades abiertas.
\bigskip


%%%%%%%%%%%%%%%%%%%%%%%%%%%%%%%%%%%%%%%%%%%%%%%%
%%                                            %%
%%           HOMOLOGIA SIMPLICIAL             %%
%%                                            %%
%%%%%%%%%%%%%%%%%%%%%%%%%%%%%%%%%%%%%%%%%%%%%%%%
\subsection{Construcción de la homología simplicial}

Una ligera generalización de los complejos simpliciales que dan nombre a esta teoría son los $\Delta$-complejos. Es sobre estos que textos como \cite{hatcher2002} definen la homología singular. Nosotros, como ya hemos mencionado, lo haremos para espacios topológicos triangulados y, más tarde, para CW-complejos.

\bigskip

\begin{definition}
    Sea el par $(X,\tau)$ un espacio topológico triangulado, con $\tau = \{(P^k_{\alpha}, f_{\alpha}): ~0\le k \le n, ~\alpha\in \Lambda_k\}$. Elegimos cualesquiera orientaciones $o^k_{\alpha}$ para cada poliedro $P^k_{\alpha}$ de la triangulación y llamamos $e^k_{\alpha}=(P^k_{\alpha},o^k_{\alpha})$ al poliedro orientado. Definimos $C^{\tau}_k(X)$ como el grupo libre abeliano en los generadores $e^k_{\alpha}$
    \begin{equation*}
        \mathbb{Z}\langle e^k_{\alpha} | \alpha \in \Lambda_k \rangle
    \end{equation*}
    Definimos también los siguientes homomorfismos de grupos, a los que llamaremos \textit{operadores de borde}, $\partial_k: C^{\tau}_k(X) \to C^{\tau}_{k-1}(X)$ dados por
    \begin{equation*}
        \partial_k(e^k_{\alpha}) = \sum_{\beta\in\Lambda_k} [e^k_{\alpha} : e^{k-1}_{\beta}]e^{k-1}_{\beta},
    \end{equation*}
    donde a los $[e^k_{\alpha} : e^{k-1}_{\beta}]$ los llamaremos \textit{números de incidencia} y vienen dados por $[e^k_{\alpha} : e^{k-1}_{\beta}] = \sum \varepsilon_i$ donde $i$ recorre todas las inclusiones $L_i: P^{k-1}_{\beta} \hookrightarrow \partial P^k_{\alpha}$ y $\varepsilon_i=\pm 1$ dependiendo de si la orientación $o^{k-1}_{\beta}$ coincide o no con la orientación inducida como borde por $L_i(P^{k-1}_{\beta})\subset P^k_{\alpha}$
\end{definition}

\begin{remark}\hfill
    \begin{enumerate}[label=(\roman*)]
        \item Definimos $C^{\tau}_k(X) = 0$ siempre que $k<0$ o $k>n$ y $\partial_k \equiv 0$ siempre que $k\le 0$ o $k>n$.
        \item Elegimos orientaciones cualesquiera para los poliedros dado que elegir una u otra se reduce a cambiar el generador $e^k_{\alpha}$ por $-e^k_{\alpha}$.
    \end{enumerate}
\end{remark}

\bigskip

Sin ningún bagage adicional el lector puede estar preguntándose cómo son los elementos de $C^{\tau}_k(X)$ y qué hacen exactamente los operadores de borde. Bien, llamamos \textit{$k$-cadenas} a los elementos de $C^{\tau}_k(X)$ , que son sumas formales de poliedros de $\tau$ con coeficientes enteros donde el signo depende de las orientaciones. Es decir, elegimos una orientación para que lleve el signo $+$ y los poliedros poliedros cuya orientación sea la otra posible llevarán el signo $-$. Es decir, son de la forma
\begin{equation*}
    \sum_{i=1}^m z_i e^k_{\alpha_i} ~~\text{ con }~~ z_i \in \mathbb{Z} ~~\text{ y }~~ \alpha_i \in \Lambda_k, ~\forall i\in \{1,\dots, m\}.
\end{equation*}
Arriba hemos descrito como actuan los operadores de borde $\partial_k$ sobre los generadores de $C^{\tau}_k(X)$, sin embargo, por linealidad, se pueden extender a homomorfismos de grupos de $C^{\tau}_k(X)$ a $C^{\tau}_{k-1}(X)$. Así,
\begin{equation*}
    \partial_k \left( \sum_{i=1}^m z_i e^k_{\alpha_i} \right) = \sum_{i=1}^m z_i \partial_k (e^k_{\alpha_i}). 
\end{equation*}
Ahora, el operador de borde lleva una $k$-cadena a la $(k-1)$-cadena que constituye su borde teniendo en cuenta la orientación, en el caso $2$-dimensional, diríamos que lleva un $2$-poliedro a la suma de sus aristas. Una propiedad imprescindible de los operadores de borde es la siguiente.

\begin{proposition}
    Los operadores de borde satisfacen $\partial_{k}\circ\partial_{k+1}=0$ para todo $k\ge 0$.
\end{proposition}

\begin{proof}
    Es posible demostrar que toda triangulación se puede subdividir para obtener una triangulación regular. Teniendo en cuenta que para una triangulación regular la notación y las cuentas se simplifican bastante, vamos a probar el resultado en este caso. Al no darse identificaciones entre $(k-1)$-símplices de la frontera de cada $k$-poliedro al triangular, se tiene que solo hay una inclusión posible para cada uno, es decir, $[e^k_{\alpha} : e^{k-1}_{\beta}]=\pm 1$ siempre.

    Ahora, basta comprobar que el enunciado se verifica para los generadores de $C^{\tau}_{k+1}(X)$. Sea $e^{k+1}_{\alpha}$ un $(k+1)$-poliedro orientado que genera $C^{\tau}_{k+1}(X)$, se tiene que
    \begin{equation*}
        \partial_{k}\circ\partial_{k+1}(e^{k+1}_{\alpha}) = \partial_{k}\left( \sum_{\beta\in\Lambda_{k+1}} [e^{k+1}_{\alpha} : e^{k}_{\beta}]e^{k}_{\beta} \right) = \sum_{\gamma\in\Lambda_{k}} \left( \sum_{\beta\in\Lambda_{k+1}} [e^{k+1}_{\alpha} : e^{k}_{\beta}][e^{k}_{\beta} : e^{k-1}_{\gamma}] \right) e^{k-1}_{\gamma}
    \end{equation*}
    
    Sea $e^{k-1}_{\gamma}$ un $(k-1)$-poliedro cualquiera incluido en $e^{k+1}_{\alpha}$. Como solo hay una inclusión posible de $e^{k-1}_{\gamma}$ en $\partial e^{k+1}_{\alpha}$ , existen exactamente dos $k$-poliedros orientados $e^{k}_{\beta}$ y $e^{k}_{\beta'}$ tales que 
    \begin{equation*}
        e^{k-1}_{\gamma}\subset e^{k}_{\beta} \subset e^{k+1}_{\alpha}~\text{ y }~ e^{k-1}_{\gamma}\subset e^{k}_{\beta'}\subset e^{k+1}_{\alpha}
    \end{equation*}
    Luego
    \begin{equation*}
        \partial_{k}\circ\partial_{k+1}(e^{k+1}_{\alpha}) = \sum_{\gamma\in\Lambda_{k}}\big( [e^{k+1}_{\alpha} : e^{k}_{\beta}][e^{k}_{\beta} : e^{k-1}_{\gamma}] + [e^{k+1}_{\alpha} : e^{k}_{\beta'}][e^{k}_{\beta'} : e^{k-1}_{\gamma}] \big) e^{k-1}_{\gamma}
    \end{equation*}

    Supongamos que $[e^{k+1}_{\alpha} : e^{k}_{\beta}]=[e^{k+1}_{\alpha} : e^{k}_{\beta'}]$, entonces $[e^{k}_{\beta} : e^{k-1}_{\gamma}]=-[e^{k}_{\beta'} : e^{k-1}_{\gamma}]$ luego los sumandos se cancelan.

    De la misma manera, si $[e^{k+1}_{\alpha} : e^{k}_{\beta}]=-[e^{k+1}_{\alpha} : e^{k}_{\beta'}]$ entonces $[e^{k}_{\beta} : e^{k-1}_{\gamma}]=[e^{k}_{\beta'} : e^{k-1}_{\gamma}]$ y se cancelan de nuevo. 

\end{proof}

\begin{remark}
    Llamaremos \textit{complejo de cadena} a cualquier secuencia de la forma
        \begin{equation*}
            \dots \overset{\partial_{k+2}}{\longrightarrow} A_{k+1} \overset{\partial_{k+1}}{\longrightarrow} A_{k} \overset{\partial_{k}}{\longrightarrow} A_{k-1} \overset{\partial_{k-1}}{\longrightarrow} \dots,
        \end{equation*}
        donde los $A_i$ sean grupos abelianos, como en nuestro caso, y se verifique esta propiedad. Dicha propiedad es equivalente a $\text{im}(\partial_{k+1})\subseteq\text{ker}(\partial_k)$. Diremos que la secuencia es \textit{exacta} si se da la igualdad.
\end{remark}

\bigskip

Extraemos del resultado que si una $k$-cadena es un borde, entonces esta no tiene borde. Imaginemos un disco $2$-dimensional, este es homeomorfo a un $2$-símplice, lo que nos da una triangulación. Tenemos tres vértices $v_0$, $v_1$ y $v_2$, tres aristas $[v_0,v_1]$, $[v_1,v_2]$ y $[v_2,v_0]$ y una cara $[v_0, v_1, v_2]$. Luego,
\begin{equation*}
    \partial_1(\partial_2([v_0, v_1, v_2])) = \partial_1([v_0,v_1]+[v_1,v_2]+[v_2,v_0]) = (v_1-v_0) + (v_2-v_1) + (v_0-v_2) = 0. 
\end{equation*}
El borde del $2$-símplice no tiene borde, $[v_0,v_1]+[v_1,v_2]+[v_2,v_0]\in \text{im}(\partial_2) \subseteq \text{ker}(\partial_1)$.

Así, podemos definir el $k$-ésimo grupo de homología de $X$ como sigue.

\begin{definition}
    Sea $(X,\tau)$ un espacio triangulado. El \textit{k-ésimo grupo de homología simplicial} de $X$ es el grupo abeliano $H^{\tau}_k(X)$ definido por
    \begin{align*}
        Z^{\tau}_k(X) &= ker(\partial_k: C^{\tau}_k(X) \to C^{\tau}_{k-1}(X)) \\
        B^{\tau}_k(X) &= im(\partial_{k+1}: C^{\tau}_{k+1}(X) \to C^{\tau}_{k}(X)) \\
        H^{\tau}_k(X) &= Z^{\tau}_k(X)/B^{\tau}_k(X)
    \end{align*}
\end{definition}

\begin{remark} Sea $X$ un espacio triangulado de dimensión $n$, observamos que

    \begin{itemize}
        \item $H^{\tau}_k(X)=0$ para todo $k>n$, ya que $C^{\tau}_k(X)=0$.
        \item $H^{\tau}_n(X)$ es un grupo abeliano libre. En efecto, como $C^{\tau}_{n+1}(X)=0$ se tiene que $B^{\tau}_n(X)=0$ luego $H^{\tau}_n(X)=Z^{\tau}_k(X)\subset C^{\tau}_n(X)$ que es libre por ser subgrupo de un grupo libre.
        \item Si $X$ es además compacto, entonces $H^{\tau}_k(X)$ es un grupo abeliano finitamente generado. Se tiene que $C^{\tau}_k(X)$ es isomorfo a $\mathbb{Z}^{n_k}$, donde $n_k=|\Lambda_k|$, luego es finitamente generado y todo subgrupo de un grupo abeliano finitamente generado lo es a su vez. Así, $Z^{\tau}_k(X)$ y $B^{\tau}_k(X)$ son finitamente generados y se sigue el resultado.
    \end{itemize}
\end{remark}

De nuevo, esta definición merece una aclaración. A los elementos de los subgrupos $Z^{\tau}_k(X)$ y $B^{\tau}_k(X)$ de $C^{\tau}_k(X)$ los llamaremos, respectivamente, \textit{ciclos} y \textit{bordes}. El primero lo reciben las $k$-cadenas cuyo borde es vacío (como en el caso del $2$-símplice) y el segundo aquellas que constituyen el borde de una suma formal de $(k+1)$-cadenas teniendo en cuenta la orientación. Los elementos de $H^{\tau}_k(X)$ son, por tanto, las clases de equivalencia de ciclos cuya diferencia es un borde.

Uno desearía que los grupos de homología no dependiesen de la triangulación escogida y fuesen iguales que los obtenidos por la vía de la homología singular y, en efecto, esto se cumple. 

\begin{theorem}
    Sea $(X,\tau)$ un espacio triangulado. Entonces $H^{\tau}_k(X)$ es isomorfo a $H_k(X)$ para todo $k>0$. Donde $H_k(X)$ denota el grupo de homología singular de $X$.
\end{theorem}

De nuevo, dejaremos que el lector encuentre la demostración en la fuente original en virtud de nuestra decisión de evitar formalizar la homología singular para nuestros propósitos. De ahora en adelante nos referiremos a los grupos de homología como $H_k(X)$, sin especificar una triangulación.

\bigskip

Terminemos el ejemplo sencillo del disco, proporcionaremos más ejemplos para el caso de los CW-complejos más adelante.Sea $D^2$ el disco, entonces, usando la triangulación ya mencionada formada por un $2$-símplice se tienen $C_0(D^2) = \mathbb{Z}\langle v_0, v_1, v_2 \rangle \cong \mathbb{Z}^3$, $C_1(D^2)=\mathbb{Z}\langle [v_0,v_1],[v_1,v_2],[v_2,v_0] \rangle \cong \mathbb{Z}^3$ y $C_2(D^2)=\mathbb{Z}\langle [v_0,v_1,v_2] \rangle \cong \mathbb{Z}$. 
    
Ahora, $Z_2(D^2)= \text{ker}(\partial_2)=\{0\}$ ya que los elementos de $C_2(D^2)$ son de la forma $z([v_0,v_1,v_2])$ con $z$ entero y $\partial_2(z[v_0,v_1,v_2]) = z([v_0,v_1]+[v_1,v_2]+[v_2,v_0]) \neq 0$ si $z$ es distinto de cero. Además, $B_2(D^2) =\{0\}$ necesariamente luego $H_2(D^2)=\{0\}$.
    
Continuamos, $B_1(D^2)=\text{im}(\partial_2)=\{z[v_0,v_1,v_2]\}_{z\in\mathbb{Z}}\cong \mathbb{Z}$ y $Z_1(D^2)= \text{ker}(\partial_1)=B_1(D^2)$ luego $H_2(D^2)\cong \mathbb{Z}/\mathbb{Z} = \{0\}$.

No hace falta comprobar que $Z_0(D^2)=\text{ker}(\partial_0)=C_0(D^2)\cong \mathbb{Z}$. Por otro lado, $B_0(D^2)=\text{im}(\partial_2)\subseteq \text{ker}(\partial_1)$ luego es necesariamente $\{0\}$. Así, $H_0(D^2)\cong \mathbb{Z}$.

¿Qué ocurre si eliminamos el interior del $2$-símplice? Obtenemos una circunferencia $\mathbb{S}^1$ triangulada por tres $1$-símplices y tres $0$-símplices. Ahora, los grupos $C_1$ y $C_0$ se mantienen, así como $Z_1$, sin embargo, $B_1(\mathbb{S}^1)$ ahora es $\{0\}$ necesariamente por lo que $H_1(\mathbb{S}^1)\cong\mathbb{Z}/\{0\} = \mathbb{Z}$. Qué coincide con el resultado esperado ya que ahora $[v_0,v_1]+[v_1,v_2]+[v_2,v_0]$ es un ciclo que no es borde de nada, luego representa una clase no trivial del grupo de homología.

\bigskip

Antes de seguir, listamos algunos resultados que podrían ser de interés para el lector en el contexto de las variedades diferenciables y la teoría de Morse. Sin embargo, decidimos no entreternos ya que no es central a nuestro objetivo y los presentamos sin demostración:

\begin{enumerate}
    \item Para cualquier espacio triangulado $X$, $H_0(X)=\mathbb{Z}^r$ donde $r$ es el número de componentes conexas por caminos de $X$.
    \item Si $M$ es una $n$-variedad compacta, conexa, orientada y sin borde triangulada, entonces $H_n(X)=\mathbb{Z}$.
    \item Si $M$ es una $n$-variedad compacta, conexa, no orientada y sin borde triangulada, entonces $H(X)=0$.
    \item Si $M$ es una $n$-variedad compacta, conexa y con borde (no vacío), entonces de nuevo $H_n(X)=0$.
\end{enumerate}

 \bigskip

%%%%%%%%%%%%%%%%%%%%%%%%%%%%%%%%%%%%%%%%%%%%%%%%
%%                                            %%
%%            HOMOLOGIA CELULAR                %%
%%                                            %%
%%%%%%%%%%%%%%%%%%%%%%%%%%%%%%%%%%%%%%%%%%%%%%%%
